\chapter{Simulation 01: three interacting points}\label{ch:s01}
\plain{Three points start at different random locations, at rest. Each point
responds to the other two, adds that response to its retained velocity,
damps the result, and moves. The points take turns. Only after all 500 rounds
does the program draw their complete trails. The trail is a record, not an
extra moving object or a route the points consult.}

\section{Entities, initialization and retained state}
The complete 74-line \filepath{recursive_attractor.py} imports NumPy and
Matplotlib, defines three helpers, performs the whole numerical loop at
module level, then makes a static plot. There is no main guard or animation
callback. Importing the file would execute the simulation; no such import
was made for this volume. \src{S01}{1--74}; \cite{construction01}.
\claim{C09}

There are three position vectors $x_i\in[-1,1)^2$ and three velocity vectors
$v_i\in\mathbb R^2$, giving twelve real-valued state entries. The source
initializes six coordinates by \filepath{np.random.rand(3,2)*2-1},
with no seed statement, and sets all six velocity entries to zero.
It sets $N=3$, $\Delta t=0.05$ and 500 sweeps.
There are no masses, radii, collision surfaces, binary modes or permanent
leader/follower assignments. The legend later labels array indices $0,1,2$
as Attractors $1,2,3$. \src{S01}{5--14,66--71}. \claim{C10}

The current arrays carry state between sweeps. The initially empty
\filepath{position_log} stores observations for drawing. Neither that log
nor its earlier entries are read by the movement rule. A point's velocity
therefore retains consequences of previous interactions, but the program
does not contain an explicit remembered itinerary.
\src{S01}{14,38--62}. \claim{C11}

\section{The exact operators}
The source's function H is
\begin{equation}\label{eq:H}
H(t)=2e^{-0.5t}\cos t+e^{-0.3t}\sin t
      +1.5e^{-0.7t}+e^{-0.4t}.
\end{equation}
Its four amplitude/decay pairs are $(2,0.5)$, $(1,0.3)$, $(1.5,0.7)$
and $(1,0.4)$. It depends only on its time argument. The name “REFU memory
curve” in the source comment does not add a dependence on recorded
positions. \src{S01}{16--27}. \claim{C12}

Define componentwise wrapping and damping by
\begin{align}
\W_2(z)&=((z+1)\bmod 2)-1,\label{eq:W2}\\
\R(u)_a&=\frac{u_a}{1+\log(1+|u_a|)}.\label{eq:RPCO}
\end{align}
Equation~\eqref{eq:W2} is TGMO, used both on relative displacements and on
new positions. Its range is half-open; an exact half-period tie selects
$-1$ in that coordinate. Equation~\eqref{eq:RPCO} acts separately on
components, not on the Euclidean norm of a whole vector.
\src{S01}{29--35}. \claim{C13}

For one sweep $n$, let $x_i^n,v_i^n$ denote the incoming state and
$x_i^{n+1},v_i^{n+1}$ the values after body $i$ has been updated. Before
updating body $i$, define the actually available positions
\begin{equation}\label{eq:mixed01}
\widetilde x_j^{\,n,i}=
\begin{cases}
x_j^{n+1}, &j<i,\\
x_j^n, &j\geq i.
\end{cases}
\end{equation}
For $i=0,1,2$ in that order, the exact mathematical transcription is
\begin{align}
t_n&=0.05n,\qquad h_n=H(t_n),\label{eq:clock01}\\
d_{ij}^{\,n,i}&=\W_2(\widetilde x_j^{\,n,i}-x_i^n),\qquad j\ne i,\label{eq:pairdiff}\\
F_i^{\,n}&=\sum_{j\ne i}
 \frac{d_{ij}^{\,n,i}}{\|d_{ij}^{\,n,i}\|_2+0.001},\label{eq:pairforce}\\
u_i^{\,n}&=v_i^n+0.05h_nF_i^{\,n},\label{eq:kick01}\\
v_i^{n+1}&=\R(u_i^{\,n}),\qquad
x_i^{n+1}=\W_2(x_i^n+0.05v_i^{n+1}).\label{eq:move01}
\end{align}
The shared $h_n$ is evaluated once per sweep. The regularizer is added
\emph{after} taking the norm; it is not a squared-distance softening.
At coincidence the corresponding displacement contribution is zero.
No inverse-square law, potential or physical unit is supplied by these
expressions. \src{S01}{38--56}. \claim{C14}

\section{One sweep in execution order}
At the beginning of a sweep, the source creates an empty
\filepath{new_positions} list. For each body it resets a two-component
interaction sum, visits the other indices, skips itself, wraps each
difference and adds its regularized direction. It then performs
velocity addition, componentwise damping and position wrapping, and appends
a \emph{copy} of the resulting position. \src{S01}{42--57}.

\begin{table}[htbp]\centering
\caption{Position dependencies during one S01 sweep. A common clock does
not imply simultaneous numerical updates.}\label{tab:order01}
\begin{tabular}{@{}lll@{}}\toprule
Body updated & First other position & Second other position\\\midrule
0 & Old body 1 & Old body 2\\
1 & New body 0 & Old body 2\\
2 & New body 0 & New body 1\\\bottomrule
\end{tabular}
\end{table}

The list is an observation buffer, not an array of proposals committed
simultaneously. Earlier mutations are already visible to later bodies.
After the third move, the three copied positions form an end-of-sweep
snapshot because later bodies do not overwrite earlier rows.
The final log has shape $(500,3,2)$ and excludes the initializer.
It stores neither velocities nor within-sweep comparisons.
\src{S01}{43--62}. \claim{C15}

The first stored state was produced with $t=0$; the last with $t=24.95$.
Counting 500 multipliers of $0.05$ gives 25 model-step units, but no
physical-time calibration. The program stops because its loop ends, not
because it detects equilibrium. A zero initial velocity is not restored
after each sweep. \src{S01}{5--11,38--59}. \claim{C16}

\section{What the historical drawing establishes}
After simulation, one plot call per body connects its 500 recorded positions
into 499 line segments. The script supplies title, legend, grid, equal
axis scaling and a final \filepath{plt.show()}. It supplies no
\filepath{savefig}, CSV export, movie writer, temporal arrows or
initial-state marker. The saving method for the surviving
\filepath{attractors.png} is missing. \src{S01}{64--74}. \claim{C17}

Figure~\ref{fig:F01} is the original attributed image. Its exact title
and three traces match the source construction. Closely aligned upper
curves, separate turns and long near-vertical central connections are
visible. Those features do not disclose the order in which the points
visited a location. Since wrapped samples are connected directly,
the source permits a seam crossing to produce a long flat-plot connector.
Without the generating numerical log, that is an explanation the source
permits, not a certified diagnosis of each line in the old PNG.
\cite{construction01} \claim{C18}
\EvidenceFig{F01}

The plotted line is also a discrete drawing convention. It is not a
continuous-time solution supplied by the code. Neither overlapping
segments nor the apparent merging of colors proves simultaneous
occupancy, exact coincidence, a bound pair or route inheritance.
The original seed and ordered history were not recovered in the bounded
construction study. That negative recovery result does not prove no
such evidence exists elsewhere. \claim{C19}

\section{The separately accepted seed-0 reproduction}
The subsequent replay study was explicitly authorized as one new
500-sweep configuration using legacy NumPy seed~0. Its initializer is
shown in Table~\ref{tab:init01}. The historical rules, zero velocities,
arithmetic order, clock and wrapping were preserved. The accepted
environment was Python~3.11.15, NumPy~2.4.4, Matplotlib~3.10.9 and
float64 arrays, with Agg rendering. \cite{replay01} \claim{C20}

\begin{table}[htbp]\centering\small
\caption{The reproduced initializer, not the missing historical initializer.
All initial velocities are $(0,0)$. Full floating-point values and hexadecimal
representations remain in replay/initial\_state.json.}\label{tab:init01}
\begin{tabular}{@{}lrr@{}}\toprule
Body & X & Y\\\midrule
0 &0.0976270078546495&0.43037873274483895\\
1 &0.20552675214328775&0.08976636599379373\\
2 &$-0.15269040132219058$&0.29178822613331223\\\bottomrule
\end{tabular}
\end{table}

The accepted harness recorded 501 sweep-boundary states including $S_0$,
1,500 ordered point events, incoming and outgoing velocities, comparisons,
native increments and integer wrap shifts. A single same-seed reference
execution of the unchanged copied source supplied a fidelity comparison:
all 3,000 logged position components and six final velocity components
matched bit for bit. Intermediate reference velocities were not compared
at every sweep. The present volume reads that receipt; it performs
neither execution again. \claim{C21}

Figure~\ref{fig:F02} is therefore labelled reproduced, not original.
Its long shared diagonal belongs to this known seed. The very different
historical picture in Figure~\ref{fig:F01} cannot be assigned this
initializer or chronology. \claim{C22}
\EvidenceFig{F02}

\subsection{The opening updates}
The first inspection window was fixed as events 1--9, the first three
sweeps. In event~1, body~0 moves by approximately
$(-0.005728,-0.012602)$. Body~1 then sees its changed coordinates and
moves by approximately $(-0.010860,+0.012595)$. Body~2 sees both updated
positions and moves by approximately $(+0.014857,-0.000239)$.
These are saved records, not new calculations for this manuscript.
Figure~\ref{fig:F03} makes the ordered dependence visible.
\cite{replay01}, within\_sweep.jsonl lines 1--9. \claim{C23}

For body~0's first update, the accepted record gives approximately
\[
\begin{aligned}
F_0&=(-0.570665,-1.433338),\\
u_0&=(-0.128400,-0.322501),\\
v_0^+&=(-0.114561,-0.252047).
\end{aligned}
\]
The damped result is the next incoming velocity.
All three points move in the opening sweep of this reproduction;
this observation is not a general theorem excluding small-motion
waiting under every initializer. \claim{C24}
\EvidenceFig{F03}

\subsection{Near-contact and subsequent separation}
The accepted replay selected the smallest simultaneous periodic
separation among the 500 recorded end-of-sweep states and three unordered
pairs, with a declared tie rule. It selected bodies~1 and~2 at $S_{169}$.
The inspected events are 502--510, with the preceding state at event~501.
The central sweep uses source step~168, $t=8.4$ and
$H(t)\simeq0.09211406$. \claim{C25}

\begin{table}[htbp]\centering
\caption{A finite reproduced near-contact. Distances are in coordinate
units at recorded sweep boundaries, not interpolated continuous minima.
Source: replay/observation\_summary.json and within\_sweep.jsonl events
504,507,510.}\label{tab:near 01}
\begin{tabular}{@{}lrr@{}}\toprule
State & Completing event & Distance between bodies 1 and 2\\\midrule
$S_{168}$ &504&0.001496078692\\
$S_{169}$ &507&0.0000442435133\\
$S_{170}$ &510&0.001583353616\\\bottomrule
\end{tabular}
\end{table}

At the central sweep, body~0's interaction points up-right but its
outgoing velocity remains approximately $(-0.040387,-0.040422)$.
Bodies~1 and~2 have down-left interaction sums while their outgoing
velocities remain approximately $(0.008924,0.008872)$ and
$(0.030815,0.030547)$. Carried velocities explain why the instantaneous
interaction direction does not identify the movement direction.
The two nearly overlapping markers conceal distinct coordinates and
velocities; the next sweep separates them again.
\cite{replay01}, events 505--507 and 510. \claim{C26}
\EvidenceFig{F04}

This is a measured near-contact in one reproduced chronology. It is not
exact coincidence, a proven continuous collision, or a delayed
leader--pair route-following episode. The point moving the other way
along the same diagonal demonstrates why a completed line picture is
insufficient for that interpretation. \claim{C27}

\subsection{The complete finite record}
The accepted record has no position seam crossings in 1,500 point
updates and no nonzero wrap shift in 3,000 interaction comparisons.
The wrapping operators still execute. The largest native increment is
about 0.03456666287 at event~20, body~1 in source step~6.
The ordinary and seam-aware paths therefore have the same curves for
this initializer. This absence is not a property proved for every
execution of S01. \claim{C28}

Figure~\ref{fig:F05} shows repeated approach and separation together with
smaller later increments. All final velocities at $S_{500}$ remain
nonzero. The accepted record also contains seven sign changes of its
sampled H values; no ablation was performed to assign a full episode
to H alone. Finite slowing does not establish eventual convergence.
\cite{replay01}, observation\_summary.json and record\_checks.json.
\claim{C29}
\EvidenceFig{F05}
\FloatBarrier
